\FloatBarrier
\section{Further results for the U.S.}\label{app:USrob}\setcounter{equation}{0}
This appendix complements sections \ref{sec:oecd} and \ref{sec:esc} for the U.S.\ model with the counterfactuals under general CRRA preferences, the alternative calibration of the taste for leisure, and the endogenous design under ageing. The counterfactuals that impose France's characteristics are in appendix \ref{app:UKUS}, next to the same exercise on the UK.

\subsection{CRRA preferences}\label{app:US:CRRA}
Table \ref{table:US:CRRA:pensChars} repeats the pension-design counterfactual of table \ref{table:US:pensChars} at $\rho = 0.5$ and $2$, each $\rho$ separately calibrated. The effect of the design on the tax rate falls with the intertemporal elasticity, for the reason section \ref{sec:oecd} gives: moving from $\theta = 1$ to $\theta = 0$ raises the tax rate by $9.0$ p.p.\ at $\rho = 0.5$, $6.2$ p.p.\ at $\rho = 1$ and $1.5$ p.p.\ at $\rho = 2$. The savings rate and the workweek follow the tax rate, and at $\rho = 2$ the savings rate barely moves in either direction.

%% GENERATED by python/paper/build.py from results/ -- do not edit by hand.
%% Rebuild:  .venv\Scripts\python.exe python\paper\build.py --only US_CRRA_PensChars
%% Source:   results/shocks/US_shocksCommonX.csv
\begin{table}[!htb]
\centering
\begin{threeparttable}
\caption{Does CRRA matter for the effect of pension design ($\theta$) in US -- 2020}
\label{table:US:CRRA:pensChars}
\renewcommand{\arraystretch}{1.25}
\begin{tabularx}{.9\textwidth}{YYYYY}
\toprule
\textbf{Scenario} & \textbf{CRRA} ($\rho$) & \textbf{Tax rate} & \textbf{Savings rate} & \textbf{Avg. workweek} \\
\midrule 
 & 0.5 & 14.43\% & 15.91\% & 39.39 \\
Baseline & 1.0 & 14.43\% & 15.37\% & 39.39 \\
 & 2.0 & 14.43\% & 14.85\% & 39.39 \\[.5em]\hline\\[-.75em]
 & 0.5 & 20.23\% & $-2.18$ p.p. & 36.80 \\
$\theta = 0$ & 1.0 & 18.52\% & $-0.93$ p.p. & 37.52 \\
 & 2.0 & 15.21\% & $+0.10$ p.p. & 38.38 \\[.5em]\hline\\[-.75em]
 & 0.5 & 11.21\% & $+1.54$ p.p. & 40.55 \\
$\theta = 1$ & 1.0 & 12.28\% & $+0.65$ p.p. & 40.13 \\
 & 2.0 & 13.72\% & $+0.06$ p.p. & 39.79 \\[.5em]\hline\\[-.75em]
\bottomrule
\end{tabularx}
\begin{tablenotes}[flushleft]
\footnotesize
\item[] \textit{Note:} Every $\rho$ is separately calibrated. Every scenario row reports the change in the savings rate against the baseline at the same $\rho$, in percentage points.
\end{tablenotes}
\end{threeparttable}
\end{table}


Table \ref{table:US:CRRA:ageing} does the same for ageing, whose effect varies little with the elasticity. Acute ageing raises the tax rate by $9.3$, $8.7$ and $7.6$ p.p.\ at $\rho = 0.5$, $1$ and $2$, and mild ageing by $4.3$, $4.0$ and $3.6$ p.p.; the savings rate falls by between $1.8$ and $2.3$ p.p.\ of GDP under acute ageing, and the workweek by at most $0.7$ hours in either scenario.

%% GENERATED by python/paper/build.py from results/ -- do not edit by hand.
%% Rebuild:  .venv\Scripts\python.exe python\paper\build.py --only US_CRRA_Ageing
%% Source:   results/shocks/US_shocksCommonX.csv
\begin{table}[!htb]
\centering
\begin{threeparttable}
\caption{Does CRRA matter for the effect of ageing in US -- 2020}
\label{table:US:CRRA:ageing}
\renewcommand{\arraystretch}{1.25}
\begin{tabularx}{.9\textwidth}{YYYYY}
\toprule
\textbf{Scenario} & \textbf{CRRA} ($\rho$) & \textbf{Tax rate} & \textbf{Savings rate} & \textbf{Avg. workweek} \\
\midrule 
 & 0.5 & 14.43\% & 15.91\% & 39.39 \\
Baseline & 1.0 & 14.43\% & 15.37\% & 39.39 \\
 & 2.0 & 14.43\% & 14.85\% & 39.39 \\[.5em]\hline\\[-.75em]
 & 0.5 & 18.76\% & $-1.10$ p.p. & 39.08 \\
Mild ageing & 1.0 & 18.45\% & $-1.00$ p.p. & 39.08 \\
 & 2.0 & 18.01\% & $-0.88$ p.p. & 39.08 \\[.5em]\hline\\[-.75em]
 & 0.5 & 23.75\% & $-2.32$ p.p. & 38.68 \\
Acute ageing & 1.0 & 23.11\% & $-2.08$ p.p. & 38.66 \\
 & 2.0 & 21.99\% & $-1.79$ p.p. & 38.72 \\[.5em]\hline\\[-.75em]
\bottomrule
\end{tabularx}
\begin{tablenotes}[flushleft]
\footnotesize
\item[] \textit{Note:} Every $\rho$ is separately calibrated. Every scenario row reports the change in the savings rate against the baseline at the same $\rho$, in percentage points.
\end{tablenotes}
\end{threeparttable}
\end{table}


\FloatBarrier
\subsection{The vector-\texorpdfstring{$X_i$}{Xi} calibration}\label{app:US:vectorX}
Section \ref{sec:oecd} takes the taste for leisure to be common across income groups and fixes its level by the average workweek, so that relative hours are a prediction. The alternative treats relative hours as data: with both $z_i^{\eta}$ and $h_i/h$ observed, \eqref{eq:calibration_eta} identifies one $X_i$ per group, and the level of hours is then not identified. Since the productivity distribution enters every aggregate only through $\sum_i\gamma_i\eta_i^{1+\xi}/X_i^{\xi}$, which both calibrations normalize to one, $\beta$, $\omega$, the tax and savings rates and the pension-design, ageing and voting counterfactuals are the same under the two, to every printed digit for the U.S.\ and to $0.01$ p.p.\ for the UK. Only the counterfactuals defined through $\eta_i$ or $X_i$ separately read differently: France's taste for leisure moves hours alone under either calibration, and France's income distribution, alone or in the composite, is the one whose effect on the tax rate changes. Under the alternative it lowers the U.S.\ tax rate by $1.6$ p.p.\ rather than $1.2$ p.p.\ at $\rho = 1$, and it lowers the UK's tax rate where under a common $X$ it raises it (figure \ref{fig:robustness}). Online Appendix \oa{oecd-vectorx} describes the alternative, and Online Appendix \oa{oecd} shows every table and figure of section \ref{sec:oecd} under both calibrations.

\FloatBarrier
\subsection{Endogenous pension design}\label{app:US:escTables}

Table \ref{table:US_ESC:ageing} reports the acute-ageing counterfactual of section \ref{sec:esc} with the design pinned at the U.S.\ value and with the design chosen, at $\rho \in \{0.5, 1, 2\}$; figure \ref{fig:US_ESC:overview} collects it with the counterfactuals that impose France's characteristics, whose tables are in appendix \ref{app:US:ukESC}. The chosen design rises from $0.738$ to $0.816$ at $\rho = 0.5$ and to $0.825$ and $0.826$ at $\rho = 1$ and $2$, so the response to ageing hardly depends on the elasticity, and the tax rate ends between $0.1$ and $0.3$ p.p.\ above the pinned reading, since a more Bismarckian system dissipates less of its revenue and is worth a larger tax to the electorate. The additional fall in the savings rate is at most $0.4$ p.p.\ of GDP, and the workweek is unchanged to a tenth of an hour.

%% GENERATED by python/paper/build.py from results/ -- do not edit by hand.
%% Rebuild:  .venv\Scripts\python.exe python\paper\build.py --only US_ESC_Ageing
%% Source:   results/esc/escExperiments.csv
\begin{table}[!htb]
\centering
\begin{threeparttable}
\caption{Endogenous design and ``acute ageing'' in the U.S.}
\label{table:US_ESC:ageing}
\renewcommand{\arraystretch}{1.25}
\begin{tabularx}{\textwidth}{p{2.6cm}YYYYY}
\toprule
\textbf{Scenario} & \textbf{CRRA} ($\rho$) & $\bm{\theta}$ \textbf{(2020)} & \textbf{Tax rate} & \textbf{Savings rate} & \textbf{Avg. workweek (hours)} \\
\midrule 
 & 0.5 & 0.738 & 14.41\% & 16.10\% & 39.45 \\ % row: baseline@0.5
Baseline & 1.0 & 0.738 & 14.43\% & 15.41\% & 39.47 \\ % row: baseline@1.0
 & 2.0 & 0.738 & 14.44\% & 14.87\% & 39.42 \\[.5em]\hline\\[-.75em] % row: baseline@2.0
 & 0.5 & 0.738 & 23.01\% & $-1.91$ p.p. & 39.00 \\ % row: pinned@0.5
Exogenous $\theta$ & 1.0 & 0.738 & 22.51\% & $-1.85$ p.p. & 38.83 \\ % row: pinned@1.0
 & 2.0 & 0.738 & 21.79\% & $-1.76$ p.p. & 38.76 \\[.5em]\hline\\[-.75em] % row: pinned@2.0
 & 0.5 & 0.816 & 23.32\% & $-2.28$ p.p. & 38.97 \\ % row: chosen@0.5
Endogenous $\theta$ & 1.0 & 0.825 & 22.80\% & $-2.04$ p.p. & 38.87 \\ % row: chosen@1.0
 & 2.0 & 0.826 & 21.92\% & $-1.81$ p.p. & 38.84 \\[.5em]\hline\\[-.75em] % row: chosen@2.0
\bottomrule
\end{tabularx}
\begin{tablenotes}[flushleft]
\footnotesize
\item[] \textit{Note:} The deadweight cost $f(\theta_t, \tau_t)$ of equation \eqref{eq:esc:budget} is quadratic in the implicit tax the flat component levies and scaled by the size of the system, $f = \exp\{-\tfrac12 \lambda \tau_t \tilde V (1-\theta_t)^2\}$, with $\lambda$ calibrated per $\rho$ (table \ref{table:US_ESC:calibration}). Each $\rho$ is separately calibrated and its workweek normalised against its own baseline. The savings rate is savings relative to GDP; the baseline rows report its level and every other row the change against the baseline at the same $\rho$, in percentage points. The ``acute ageing'' scenario sets $\nu_t = 1$ throughout.
\end{tablenotes}
\end{threeparttable}
\end{table}


A cost attached to the design itself rather than to the transfer it performs, $f(\theta) = \phi + (1-\phi)\theta^{p}$, is the alternative specification of Online Appendix \oa{esc-scale}; figure \ref{fig:robustness} sets the designs chosen under it beside those of the cost of section \ref{sec:esc:cost}.
